
\section{Design}


\label{sec:design}
In this section, we present the two-tier compiler-serving co-design of \sysname. As shown in Figure~\ref{fig:design_overview}, the \emph{agent descriptor} is generated by combining analysis of the agent program with observation of runtime behaviors. The server side then exploits the derived descriptor to schedule the KV-cache efficiently, using \emph{adaptive prompt re-layout} (\S\ref{sec:relayout}) and \emph{proactive KV-cache scheduling} (\S\ref{sec:Proactive KV-Cache Scheduling}).


\subsection{Agent Descriptor}
The \emph{agent descriptor} contains the metadata for scheduling. Static analysis provides its skeleton, and the server refines the weights with runtime observations. An agent descriptor contains two parts: 
\begin{enumerate*}[label=(\roman*)]
    \item \textbf{Call-site Template:} It provides a heuristic for each component's lifetime and sharing pattern.
    \item \textbf{Workflow Automaton:} It contains the program-level structure of an agent, including workflow dependencies, execution time estimation and transition probabilities.
\end{enumerate*} 

\subsubsection{Call-site Template}
\label{sec:template}
Prompts can be decomposed into heterogeneous components. We classify these components into the following two categories:
\begin{enumerate*}[label=(\roman*)]
    \item \textbf{Invariant Specifications:} the part of prompt that is fixed once the program is compiled. They are rendered from the compiler's own intermediate representation and stay constant.
    \item \textbf{Dynamic Bindings:} the argument values of LLM functions, which are determined at runtime.
    % \item \textbf{Dynamic Bindings:} because LLM calls are declared as ordinary functions, the arguments of a call-site are supplied at runtime. These runtime values are the dynamic bindings.
\end{enumerate*}

We define the call-site template as $\mathcal{T}_c=\langle \mathcal{I}_c,\ \mathcal{B}_c,\ \mathcal{H}_c \rangle$, where $\mathcal{I}_c$ denotes the invariant specifications, $\mathcal{B}_c$ the dynamic bindings, and $\mathcal{H}_c$ the heterogeneity annotation over those bindings.

\paragraph{Heterogeneity annotation}
For each dynamic binding, $\mathcal{H}_c$ records both how its value is shared and how long that value stays valid: $\mathcal{H}_c:\ \mathcal{B}_c \longrightarrow \mathcal{E} \times \mathcal{S}$, where $\mathcal{E}$ is the evolution component that governs which bytes a later prompt can reuse, and $\mathcal{S}$ is the scope of this binding in code, which decides when cached bytes can be given up.

We divide $\mathcal{E}$ into five classes, listed from the most to the least reusable:
\begin{enumerate*}[label=\textbf{(\roman*)}]
    \item \textbf{CONST:} the value is fixed by the program text and is identical at every call.
    \item \textbf{COPY:} the value is the same as the value another call-site carries.
    \item \textbf{RESP:} the value comes from the reply produced by another call-site.
    \item \textbf{EXTEND:} the binding is an append-only context, where the previous value remains a prefix of the new one.
    \item \textbf{VOLATILE:} the value is freshly produced at each call, or its origin lies outside the analysis boundary.
\end{enumerate*}

% We write $\rho$ for the rank of a class in the order above, which expresses how much of a binding's bytes an earlier call can already supply. 
\subsubsection{Workflow Automaton}
\label{sec:TPFM}
We abstract an agent's workflow as a time-annotated probabilistic finite-state automaton. We formulate it as: 
\[
\mathcal{A} = (Q, P, q_0, F, \mathcal{T}),
\]
where:
\begin{enumerate*}[label=(\roman*)]
    \item $Q$ is the set of states. Each state $q\in Q$ is labeled by an LLM call-site;
    \item $P$ denotes the transition probabilities between states;
    \item $q_0 \in Q$ is the entrance call-site;
    \item $F$ specifies the termination probability of the agent at each state;
    \item $\mathcal{T} = \left\{D^{\mathrm{exec}}_q, D^{\mathrm{gap}}_{q,q'}\right\}$ is the timing annotations associated with states and transitions, where $D^{\mathrm{exec}}_q$ is the execution-time distribution of the LLM call-site $q$, and $D^{\mathrm{gap}}_{q,q'}$ is the inter-call gap distribution between the completion of $q$ and the invocation of its successor $q'$.
\end{enumerate*}

% Static analysis fixes the shape of this automaton, while its weights come from runtime logs. The server observes $\mathcal{T}$ and $P$ online and accumulates them across sessions, updating the statistical model of the agent consistently. 

\subsection{Agent Descriptor Skeleton via Compiler Analysis}
\label{sec:analysis}
% The skeleton of the agent descriptor comes from program code: the graph structure of $\mathcal{A}$ and the tuple $\mathcal{T}_c$ for each call-site $c$.

\para{Call-site graph construction}
The compiler identifies LLM call-sites from the typed AST. It then builds the interprocedural control-flow graph (CFG) of the agent and projects it onto these call-sites. An edge $(q, q')$ is added whenever $q'$ is directly reachable in one hop from $q$, while non-LLM function call-sites are contracted. If a call target cannot be resolved statically, the compiler conservatively adds edges to all candidate targets. This projected graph forms the unweighted structure of $\mathcal{A}$. Finally, the compiler applies Bourdoncle's weak topological ordering~\cite{bourdoncle1993efficient} to recover the loop nesting structure from the cycles in $\mathcal{A}$.

\para{Separating $\mathcal{I}_c$ and $\mathcal{B}_c$}
$\mathcal{I}_c$ and $\mathcal{B}_c$ are obtained by running the lowering process that generates actual prompts. The compiler runs the framework's lowering pass ahead of time using symbolic placeholders for the arguments to produce the static portion of the runtime prompt ($\mathcal{I}_c$) while leaving one labeled placeholder per parameter to represent $\mathcal{B}_c$.

\para{Recovering evolution classes ($\mathcal{E}$)}
For each binding $b \in \mathcal{B}_c$, the compiler computes the reaching definitions of its underlying variable across the CFG (including back edges, with local definitions shadowing earlier ones). Each reaching definition is classified by its right-hand side:
\begin{enumerate*}[label=(\roman*)]
    \item literals and module-level constants are classified as \textbf{CONST};
    \item variables passed unmodified from an earlier call-site $c'$ are classified as \textbf{COPY}($c'$);
    \item values assigned from the return value of another call-site $c'$ (traced via def--use chains) are classified as \textbf{RESP}($c'$);
    \item prefix-preserving updates, like string/list concatenation ($v \leftarrow v + \Delta$), are classified as \textbf{EXTEND};
    \item all other definitions are defaulted to \textbf{VOLATILE}, such as external input.
\end{enumerate*}

If different CFG paths yield conflicting evolution classes for $b$, the compiler falls back to \textbf{VOLATILE}. To let the runtime resolve this conservativeness once the actual transition $(q, q')$ is known, the compiler attaches transition-specific metadata to each edge $(q, q')$. Specifically, it records which and how the variables are modified along $(q,q')$, and attaches the exact evolution class for each binding at $q'$. For example, a binding can be reduced to \textbf{CONST} rather than \textbf{VOLATILE} if the exact edge assigns it a literal.


\para{Recovering lifetime scopes ($\mathcal{S}$)}
The compiler records $\mathcal{S}(b)$ as the set of procedures (or loop bodies) that overwrite the binding $b$'s underlying variable. When the variable is assigned once and remains valid for the entire session, $\mathcal{S}_(b)=\emptyset$. Whenever execution enters a procedure in $\mathcal{S}(b)$, the variable is overwritten and its previously cached KV bytes become dead.

\subsection{Adaptive Prompt Re-layout}
\label{sec:relayout}
\begin{algorithm}[t]
    \caption{Layout decision algorithm}
    \label{alg:layout}
    \begin{algorithmic}[1]
    \algtext*{EndFor}\algtext*{EndIf}\algtext*{EndWhile}
    \Require Descriptor $(\mathcal{A},\{\mathcal{T}_q\}_{q\in Q})$, switch \textit{pin}, lengths $\ell$
    \Ensure Layouts $\{\mathcal{L}_q\}_{q\in Q}$
    \State $B_{\mathrm{sh}} \gets \{b : |\{q\in Q : b\in\mathcal{B}_q\}|>1\}$\label{alg:layout:sh-begin}\Comment{shared bindings}
    \State $G \gets(B_{\mathrm{sh}},\emptyset)$ \Comment{precedence graph}
    \ForAll{$(q,q')\in\Call{Edges}{\mathcal{A}}$}
        \State $P_1,P_2,P_3 \gets \emptyset$
        \Statex \hfill Groups: unchanged, append, rewritten/one-sided
        \ForAll{$b \in B_{\mathrm{sh}}\cap(\mathcal{B}_q\cup\mathcal{B}_{q'})$}
            \State $P_{\tau(b)} \gets P_{\tau(b)}\cup\{b\}$ \Comment{$\tau(b)$: group of $b$}
        \EndFor
        \State $E(G) \gets E(G)\cup\bigcup_{i<j}P_i\times P_j$
    \EndFor
    \State $\mathcal{C} \gets \Call{Tarjan}{G}$ \Comment{SCCs of $G$, sets of bindings}
    \State $\hat{E} \gets \{(C,C')\in\mathcal{C}^2 : C\neq C',\ (C\times C')\cap E(G)\neq\emptyset\}$
    \State $\hat{G} \gets (\mathcal{C},\hat{E})$ \Comment{condensation graph of $G$}
    \State $j(b) \gets (-\ell(b),\textit{depth}(b))$ for $b\in B_{\mathrm{sh}}$
    \State $J(C) \gets \min_{b\in C} j(b)$ for $C\in V(\hat{G})$
    \State $L \gets [\,]$
    \While{$V(\hat{G})\neq\emptyset$}
        \State $C^{*} \gets \arg\min_{C\in V(\hat{G}),\ \textit{indeg}(C)=0} J(C)$
        \State $L \gets L\,\Vert\,\Call{Sort}{C^{*},j}$;\ \ $\hat{G} \gets \hat{G}-C^{*}$ \label{alg:layout:sh-end}
    \EndWhile
    \ForAll{$q\in Q$}
        \State $\pi_q \gets [\,b\in L : b\in\mathcal{B}_q\,]\,\Vert\,\Call{Sort}{\mathcal{B}_q\setminus B_{\mathrm{sh}},K}$ \label{alg:layout:pv}
        \State $s_q \gets \Call{Spec}{\mathcal{I}_q}$ \Comment{instruction specification of $q$} \label{alg:layout:spec-begin}
        \If{$\neg\textit{pin}\land\mathcal{B}_q\cap B_{\mathrm{sh}}\neq\emptyset$}
            \State $\mathcal{L}_q \gets \pi_q\Vert s_q$
        \Else
            \State $\mathcal{L}_q \gets s_q\Vert\pi_q$ \label{alg:layout:spec-end}
        \EndIf
    \EndFor
    \State \Return $\{\mathcal{L}_q\}_{q\in Q}$
    \end{algorithmic}
\end{algorithm}

% The message a framework emits for a call-site follows the order of parameters in the function declaration, which may place a changing value ahead of values that later calls could reuse, reducing the effectiveness of prefix caching. \sysname's adaptive re-layout permutes prompt components to mitigate this inefficiency.
% \sysname's \emph{adaptive prompt re-layout} permutes prompt components to maximize prefix cache reuse. 
\subsubsection{Movable Components}
First, the dynamic bindings $\mathcal{B}_c$ can be permuted freely. The compiler lowers each argument into a self-contained component labeled by its parameter name, and the framework places these components in declaration order only because that is the order of the argument list; re-layout may choose any order of the bindings.

Second, the instruction specification in $\mathcal{I}_c$ is movable. $\mathcal{I}_c$ consists of the system prompt, the response schema, the tool schema, and the instruction specification (such as task descriptions and parameter descriptions). The system prompt stays universally constant and should be placed at the very beginning, and the schemas are handled by the constrained decoding engine. The instruction specification can be put just after the system prompt or at the end of the prompt.

\subsubsection{Layout Decision}
\label{sec:layout_decision}

\sysname decides the layout $\mathcal{L}_q$ for each call-site $q$ by ranking the reusability of bindings (Algorithm~\ref{alg:layout}). It places shared bindings ahead of private bindings and prioritizes stable bindings within both categories to maximize prefix cache efficiency. 

\paragraph{Sorting shared bindings {\normalfont(lines~\ref{alg:layout:sh-begin}--\ref{alg:layout:sh-end})}} Shared bindings are those that are carried by multiple call-sites. For every transition from $q$ to $q'$, we identify how each shared binding evolves. We categorize their behaviors into three sharing pattern groups, ordered as follows:
\begin{enumerate*}[label=(\textbf{\roman*})]
    \item \textbf{unchanged:} Both call-sites carry it and its value does not change between the two calls.
    \item \textbf{append:} Both call-sites carry it and the program only appends to its value.
    \item \textbf{rewritten/one-sided:} The program rewrites it between the two calls, or only one call-site carries it.
\end{enumerate*}

\sysname builds the precedence graph of all shared bindings based on the groups they belong to on each transition edge, condenses it into a DAG using Tarjan's algorithm, and finally performs a topological sort on the DAG. We key each SCC by the smallest $\textbf{(-\textit{length}, \textit{depth})}$ among its bindings, where:
\begin{enumerate*}[label=\textbf{(\roman*)}]
    \item \textit{depth}: the number of transitions from the entrance to the first call-site that carries $b$;
    \item \textit{length}: the token length of the value of $b$, observed by the server.
\end{enumerate*}
The SCC with the smallest key goes first when multiple SCCs have zero in-degree. Each SCC is then expanded into the list of its bindings sorted also by $\textbf{(-\textit{length}, \textit{depth})}$.

\paragraph{Sorting private bindings {\normalfont(line~\ref{alg:layout:pv})}} The bindings carried by a single call-site are sorted in ascending lexicographic order of:
\[
K(b)=\big(\textit{scoped},\ \rho,\ \textit{depth},\ -\textit{length}\big),
\]
where:
\begin{enumerate*}[label=\textbf{(\roman*)}]
    \item \textit{scoped}: the lifetime of $b$; 0 if it is session-wide ($\mathcal{S}=\emptyset$) and 1 otherwise;
    \item $\rho$: the stability rank for $\mathcal{E}$ in \S\ref{sec:template}, from 1 for CONST to 5 for VOLATILE.
\end{enumerate*}


\paragraph{Placing the specification {\normalfont(lines~\ref{alg:layout:spec-begin}--\ref{alg:layout:spec-end})}} 
Placing the \emph{instruction specification} at the end gives a longer common prefix. \sysname places the specification at the end for each call-site that carries a shared binding. A switch keeps the specification at the front when the developer requires it. 



\subsection{Proactive KV-cache Scheduling}
\label{sec:Proactive KV-Cache Scheduling}
% \sysname's KV-cache scheduling proactively arranges the placement of KV-cache entries. Using the agent descriptor, the server decides operations for every served KV-cache sequence.

\subsubsection{Reusable Head and Volatile Tail}
\label{sec:identify_reusable_kv}
A served prompt naturally divides into two segments: a reusable head and a volatile tail. The reusable head consists of two parts: a \emph{static prefix} and a \emph{session-shared extension}:
\begin{enumerate*}[label=(\roman*)]
    \item \textbf{Static Prefix} is the prefix that does not change, including $\mathcal{I}_c$ and the CONST bindings that re-layout has placed behind $\mathcal{I}_c$.
    \item \textbf{Session-shared Extension} consists of the bindings after the static prefix whose values recur unchanged in a later call of the same session.
\end{enumerate*}

\subsubsection{Next-Call Prediction}
\sysname predicts three aspects of the next call:
\begin{enumerate*}[label=(\roman*)]
    \item The probability of the arrival.
    \item When the call will arrive.
    \item Which prefix of its prompt can be speculatively constructed.
\end{enumerate*}

\paragraph{Arrival probability}
Each call-site $q$ can be located by a chain of loops enclosing it; for each call-site's invocation, we mark it with a tuple $\zeta=\left(i_1, i_2, \dots, i_m\right)$, where $i_k, k \in \left[1, m\right]$ is the current iteration index of the $k$-th enclosing loop, from outermost to innermost. 

\sysname estimates the probability of the next call by a loop-tagged table. We define the table as $\mathcal{P}(q' \mid q, \zeta)$ and $\mathcal{F}(q, \zeta)$, where $\mathcal{P}$ is the probability that call-site $q$ in context $\zeta$ is followed by call-site $q'$; $\mathcal{F}$ stands for the probability that the session ends at call-site $q$ under context $\zeta$. With an empty $\zeta$ (no loop context), $\mathcal{P}$ and $\mathcal{F}$ reduce to the base transition and termination probabilities $P$ and $F$ of the automaton $\mathcal{A}$.

The table is approximated and updated by frequency counts observed at runtime. Let $q$ and $\zeta$ denote the call-site and loop context of the preceding call in the same session. When a call at $q'$ arrives, the server increments the count of outcome $q'$ in row $(q,\zeta)$; when the session closes at $q$, it increments the termination count in row $(q,\zeta)$. Each observation is also counted in every row $(q,\zeta_k)$ with $\zeta_k=(i_{m-k+1},\dots,i_m)$, $k=0,\dots,m$.

After a call at $q$ lands, we compute the probability of each outcome $o$ by Witten--Bell interpolation from the empty context $\zeta_0$ to the full context $\zeta_m=\zeta$. The outcome $o$ ranges over $\mathrm{succ}(q)$, the static successors of $q$ plus a virtual outcome $\mathrm{end}$ for session termination:

\begin{align*}
\hat{P}_k(o) &= \lambda_k\frac{c_k(o)}{n_k}+(1-\lambda_k)\hat{P}_{k-1}(o),\quad k=0,\dots,m, \\
\hat{P}_{-1}(o) &= \frac{1}{|\mathrm{succ}(q)|},\qquad \lambda_k=\frac{n_k}{n_k+d_k}\ (\lambda_k=0 \text{ if } n_k=0)
\end{align*}
$c_k(o)$ is how many times outcome $o$ followed a call at $q$ in context $\zeta_k$, $n_k=\sum_o c_k(o)$ is the total number of observations in that row, and $d_k$ is the number of distinct outcomes with $c_k(o)>0$. Then $\mathcal{P}(q'\mid q,\zeta)=\hat{P}_m(q')$ and $\mathcal{F}(q,\zeta)=\hat{P}_m(\mathrm{end})$. A context with no observation is skipped. The weight $\lambda_k$ grows with $n_k$, so a frequently observed context dominates the prediction, while a rarely observed one defers to its shorter contexts.

\paragraph{Arrival time}
The time of arrival is conservatively approximated based on recorded statistics. The workflow automaton records $D^{\mathrm{gap}}_{q,q'}$ per transition: each time a call at $q'$ arrives, the server appends the interval between the reply of the preceding call at $q$ and this arrival. The interval is the agent's own time spent in its code between two calls. We therefore estimate the time gap to the next call as:
\[
\hat{g}= \max \{t_{50}-\tfrac{1}{2}(t_{90}-t_{50}), 0 \}
\]
where $t_{50}$ and $t_{90}$ are the 50th and 90th percentiles of the gap samples on the transition $(q,q')$; The second term ensures that a transition with a more dispersed gap distribution receives a tighter estimate. 

% the predicted call is then expected to arrive $\hat{g}$ after the reply of the current call. 

\paragraph{Prefix reconstruction}
\sysname constructs the prefix of call-site $q'$ for matching KV-cache entries to schedule. Since $\mathcal{H}_c$ records the relationships among arguments, \sysname builds the prefix in layout $L_{q'}$ using the argument values observed in the current session. The construction stops at the first argument with an unknown value.


\subsubsection{KV-cache Eviction Using Lifecycle Bands}
\paragraph{Bands}
We categorize KV-cache entries into four bands, from first evicted to last:
\begin{enumerate*}[label=(\arabic*)]
    \item \textbf{RETIRED:} the volatile tails of prompts belonging to completed sessions. 
    \item \textbf{TRANSIENT:} the volatile tail of a served prompt of a live session. 
    \item \textbf{UNCLASSIFIED:} entries that analysis cannot categorize. They keep the engine's default treatment.
    \item \textbf{PLANNED:} prefixes that a predicted call will reuse. 
\end{enumerate*}



\paragraph{Band transitions}
Bands are assigned at three moments of a session's lifecycle.
\begin{enumerate*}[label=(\arabic*)]
    \item \textbf{When a served prompt enters the cache}, \sysname splits it at its \emph{reusable head} (\S\ref{sec:identify_reusable_kv}). The volatile tail and the reply generated after it become \emph{TRANSIENT}; the head stays \emph{UNCLASSIFIED} until the next prediction refresh finds a successor call that starts with it, which tries to promote it to \emph{PLANNED}.
    \item \textbf{When a later call arrives}, the transition it takes may rewrite bindings that earlier heads carried. Each affected head is cut back to the prefix that survives the rewrite, and the bytes past it are demoted to \emph{TRANSIENT}.
    \item \textbf{When the session ends}, the volatile tails of every prompt it served drop to \emph{RETIRED}. Their heads stay \emph{UNCLASSIFIED} and age out by recency, since other sessions of the same program may share them.
\end{enumerate*}

\paragraph{Band-ordered eviction}
Eviction empties the lowest band first and falls back to recency within the three lower bands. The radix tree's eviction heap is keyed by band, then by recency. The host tier keeps LRU eviction, except that \emph{RETIRED} entries go first. Additionally, KV-cache entries marked as \emph{TRANSIENT} are not admitted into the host tier as they will not be reused.

\subsubsection{KV-cache Retention and Prefetch}
% \sysname protects the required KV-cache entries based on prediction. Each prediction is summarized into a tuple of reconstructed prefix, arrival probability and estimated deadline:
A future reuse prediction is summarized as the following triple:
\[
\textit{Prediction}=\left(\pi_{q'},\ \mathcal{P}(q'\mid q,\zeta),\ t_{\mathrm{reply}}+\hat{g}\right)
\]
where $\pi_{q'}$ is the prefix reconstructed for $q'$ and $t_{\mathrm{reply}}$ is the time the current call replied. Predictions are refreshed at every call, and are dropped when the predicted calls arrive or their predicted arrival time passes. 

\sysname acts on this triple in two ways: \emph{Retention} keeps $\pi_{q'}$ from leaving the device, and \emph{prefetch} restores the part of $\pi_{q'}$ that has already been demoted to the host.

\paragraph{Retention}
\sysname ranks the predicted reusable KV-cache entries by the benefit of keeping $\pi_{q'}$ resident at every refresh:
\[
\mathcal{D}(q')=\frac{\mathcal{P}(q'\mid q,\zeta)\cdot |\pi_{q'}|_{\mathrm{cached}}}{\max\!\left(t_{\mathrm{reply}}+\hat{g}-t_{\mathrm{now}},\ \tau_{\min}\right)}
\]
where $|\pi_{q'}|_{\mathrm{cached}}$ is the number of cached tokens of $\pi_{q'}$, and $\tau_{\min}$ is the floor that prevents an imminent call from receiving an unbounded $\mathcal{D}$. \sysname admits predictions into \emph{PLANNED} in descending order of $\mathcal{D}$ until they cover a fixed fraction of the device tier. Within \emph{PLANNED}, each radix-cache node accumulates the $\mathcal{D}$ of every planned prefix that covers it, so a shared head outranks a private one.

\paragraph{Prefetch}
When part of $\pi_{q'}$ has been demoted to the host, the planner prefetches it back to the device ahead of the estimated arrival time. The prefetch release deadline is: 
\[
t_{\mathrm{release}}=t_{\mathrm{reply}}+\hat{g}-\frac{|\pi_{q'}|_{\mathrm{host}}}{B_{\mathrm{h2d}}}-\frac{|\pi_{q'}|_{\mathrm{uncached}}}{B_{\mathrm{prefill}}}
\]
where $B_{\mathrm{h2d}}$ and $B_{\mathrm{prefill}}$ are the measured host-to-device copy rate and prefill rate, $|\pi_{q'}|_{\mathrm{host}}$ is the number of tokens of $\pi_{q'}$ cached only on the host, and $|\pi_{q'}|_{\mathrm{uncached}}$ is the number of uncached tokens that require prefill.


Among released predictions, the one with earliest deadline goes first, with ties broken by the expected saving $\mathcal{P}(q'\mid q,\zeta)\cdot|\pi_{q'}|_{\mathrm{cached}}$. A prefetch is released when the engine is idle or its deadline $t_{\mathrm{release}}$ approaches. It takes only free device memory, after reclaiming \emph{RETIRED} and \emph{TRANSIENT} entries. When free space is insufficient, prefetching pauses but keeps refreshing the recency of these cache entries so that they do not age out. 
